%% RevTeX 4-2 drafting format only.
%% Replace this class with the exact editor-supplied Springer proceedings class
%% (for example, svproc, llncs, or a volume-specific class) before submission.
\documentclass[aps,pra,reprint,superscriptaddress,amsmath,amssymb,floatfix]{revtex4-2}

\usepackage{graphicx}
\usepackage{xcolor}
\usepackage{amsthm}
\usepackage[expansion=false]{microtype}
\usepackage{hyperref}
\usepackage{orcidlink}

\hypersetup{
    colorlinks=true,
    linkcolor=blue,
    filecolor=magenta,
    urlcolor=blue,
    citecolor=blue,
    pdftitle={No Trading Strategy Can Win on Every Price Path: Computability, Randomness, and the Limits of Universal Trading},
    pdfauthor={Karl Svozil},
    pdfsubject={Computational, probabilistic, and financial limits of universal trading},
    pdfkeywords={trading, no-arbitrage, diagonalization, Halting Problem, rule inference, Martin-Lof randomness, algorithmic information theory, Busy Beaver, time reversal, market impact}
}

%% Mathematical Environments
\theoremstyle{plain}
\newtheorem{theorem}{Theorem}
\newtheorem{lemma}[theorem]{Lemma}
\newtheorem{corollary}[theorem]{Corollary}
\newtheorem{proposition}[theorem]{Proposition}
\theoremstyle{definition}
\newtheorem{definition}{Definition}
\newtheorem{criterion}{Criterion}
\theoremstyle{remark}
\newtheorem{remark}{Remark}

\begin{document}

\title{No Trading Strategy Can Win on Every Price Path:\\
Computability, Randomness, and the Limits of Universal Trading}

\author{Karl Svozil\,\orcidlink{0000-0001-6554-2802}}
\affiliation{Institute for Theoretical Physics, TU Wien, Wiedner Hauptstra{\ss}e 8-10/136, A-1040 Vienna, Austria}
\email{karl.svozil@tuwien.ac.at}
\homepage{http://tph.tuwien.ac.at/~svozil}

\date{July 28, 2026}

\begin{abstract}
Every universal-trading claim pairs a trader with a market---a path, generator,
or law. For any total deterministic computable trader, its code yields a fixed
computable countermarket with proportional price moves opposing its positions;
hence no such trader wins on every computable path. For Turing-universal
market generators, a Halting-Problem reduction rules out
uniform decisions about unbounded future events, and Busy-Beaver growth rules
out a computable waiting horizon. A passive Martin-L\"of-random record is
incompressible and prevents effective test-capital processes from becoming
unbounded; no-arbitrage supplies a distinct financial boundary. Repeatable
success therefore requires a restricted market class, benchmark, risk premium,
or informational advantage. Time reversal is a simple stress test.
\end{abstract}

\maketitle

%%------------------------------------------------------------------------------
\section{Introduction}
\label{sec:intro}
%%------------------------------------------------------------------------------

Throughout this paper the elementary configuration has two sides: a
\emph{trader} and a \emph{market}. This pair is the paper's leitmotif. At each
date the trader converts the observed price history into a position; the
market supplies the next price, and their interaction determines the gain.
Depending on the question, the market side may be a realized path, a program
generating a path, a probability law over paths, or a price process constrained
by no-arbitrage. Holding the trader fixed while changing what is admitted on
the market side reveals which kind of universal claim is being made. This
organizing picture is not a claim that every financial model is literally a
two-person zero-sum game.

A strategy is called \emph{universal} here only if, after discounting by a
specified numeraire, it produces a strictly positive terminal gain from zero
initial capital on every admissible market trajectory. Zero cost,
self-financing, admissibility, and discounted gain are essential: positive
wealth obtained from a positive initial investment is not an arbitrage and
need not outperform its financing benchmark. Self-financing means that no cash
is injected later; discounting compares wealth with the funding account; and
admissibility rules out doubling systems backed by unlimited credit. The
universal claim has the quantifier order
\begin{equation}
  \exists\,\sigma\ \forall\,\mathcal E:
  \qquad \sigma\ \hbox{wins against the market }\mathcal E ,
  \label{eq:universalquantifiers}
\end{equation}
whereas the diagonal countermarket establishes
\begin{equation}
  \forall\,\sigma\ \exists\,\mathcal E_\sigma:
  \qquad \sigma\ \hbox{does not win against }\mathcal E_\sigma .
  \label{eq:counterquantifiers}
\end{equation}

The second formula does not require a market that watches orders and reacts in
real time. Fix a total deterministic computable trader. Its code can be
incorporated into a new total program that recursively generates a price path:
on the history generated so far, simulate the trader's next position and choose
the next bounded proportional price move with the opposite sign. Once the
trader is fixed, this countermarket program and its entire infinite path are
fixed as well; any finite prefix can be computed before deployment. A trader
that never takes a position---that is, does not trade at all---earns zero
rather than a strict loss, and zero already refutes a guarantee of strictly
positive gain.

This is the trading version of the standard diagonal argument against
universal computable prediction. Given a total next-bit predictor \(F\), define
the computable sequence \(x_t=1-F(x_0,\ldots,x_{t-1})\). It disagrees with
\(F\) at every step. For trading, complementing a bit becomes choosing a price
move opposite to the position. The construction is tailored to the proposed
universal rule: it proves \(\forall\sigma\,\exists\mathcal E_\sigma\), not that
one exceptional market defeats all conceivable traders.

Universal computation creates a second, distinct obstruction. If admissible
market generators can simulate arbitrary Turing machines, a future price event
can encode whether a computation halts. A program that decided every such
event would decide the Halting Problem. The Busy Beaver function sharpens the
intuition: even among programs of bounded description size, there is no
computable uniform horizon after which a still-unseen event can be certified
never to occur. Gold's
identification-in-the-limit framework and later work on unsolvable inductive
inference place this explicit reduction in a broader learning-theoretic setting
\cite{go-67,ad-91,angluin:83,2016-AlexanderKarlSvozil-ml,moore}.

A different case removes strategy-specific opposition: a binary market record
is fixed independently of the trader and assessed with respect to a specified
computable measure. If the record is Martin-L\"of random, it passes every
effective null test; equivalently, no lower-semicomputable nonnegative
supermartingale becomes unbounded on it. Algorithmic information theory gives
the complementary description that its finite prefixes cannot be compressed
by an unbounded number of bits. These are infinite-sequence statements, not
certificates that a finite sample is ``truly random,'' and they do not exclude
isolated correct forecasts or finite lucky profits. No-arbitrage supplies yet
another boundary, formulated almost surely under a pricing measure rather than
through an adversarial countermarket or an effective null test.

The market may also cease to be exogenous when the trader is large, crowded,
or publicly known: orders consume liquidity, move quotes, reveal information,
and attract imitation or opposition. Any defensible claim of repeatable
success must therefore identify a restriction, asymmetry, benchmark, or
distributional assumption absent from the unrestricted market class. This is
the operational meaning of \emph{insight} here: the trader commits to some
structure that the universal market side does not grant for free. A trend
follower posits persistence, a contrarian posits mean reversion, an option
seller relies on a risk premium, and an informed trader relies on information
not yet fully reflected in prices. The formal results do not deny finite luck;
they deny an insight-free guarantee over the unrestricted class.

\begin{table*}[t]
\caption{\label{tab:claimmap}Distinct claims organized by the trader--market
configuration. They are complementary, not interchangeable.}
\centering
\begin{ruledtabular}
\begin{tabular}{p{0.18\textwidth}p{0.27\textwidth}p{0.24\textwidth}p{0.23\textwidth}}
\textbf{Mechanism} & \textbf{Market side} & \textbf{Formal scope}
& \textbf{Conclusion} \\
\colrule
Diagonal\newline\mbox{countermarket}
& Fixed total program constructed from the fixed trader's code
& Pathwise worst case:
\(\forall\sigma\,\exists P_\sigma\)
& No total deterministic computable trader wins on every computable price path \\
Halting reduction and\newline Busy Beaver
& Total program or generator, under an explicit promise
& Individual encoded computation over an unbounded horizon
& No general computable profit or market-event decider \\
Algorithmic\newline randomness
& Passive binary record under the computable fair-coin reference law
& Each Martin-L\"of-random path; such paths have measure one
& No constructive supermartingale, hence no computable fair-game martingale,
becomes unbounded \\
No-arbitrage
& Stochastic price process admitting a coherent pricing measure
& Almost-sure payoff conditions under an ELMM
& No admissible zero-cost classical arbitrage \\
Cover universality
& Passive bounded sequence and a stated benchmark class
& Pathwise, but benchmark-relative and asymptotic
& Vanishing per-period regret, not guaranteed positive profit \\
\end{tabular}
\end{ruledtabular}
\end{table*}

The paper's principal contribution is this taxonomy and synthesis, not a claim
that the diagonal countermarket is an equilibrium model or a newly discovered
impossibility theorem. The mathematical study of unpredictable prices
runs from Bachelier through Samuelson to the Efficient Market Hypothesis
\cite{Bachelier1900,Samuelson1965,Fama1970}; the present purpose is to keep
separate several modern meanings of ``no universal strategy.'' Section
\ref{sec:computability} gives the diagonal countermarket, two Halting reductions,
the Busy-Beaver horizon, and the Martin-L\"of, computable, and Schnorr
randomness hierarchy together with its algorithmic-information interpretation.
Sections~\ref{sec:noarb} and \ref{sec:nfl} give the financial and combinatorial
boundaries. Section~\ref{sec:criterion} treats time reversal only as a practical
robustness diagnostic. The remaining sections compress the social-choice
analogy, analyze the Wheel, distinguish Cover's benchmark-relative result, and
state the conditional scope of FAPP strategies.

%%------------------------------------------------------------------------------
\section{The Trader and the Market: Computability and Randomness}
\label{sec:computability}
%%------------------------------------------------------------------------------

We now formalize the trader--market pair. The trader is a program
\(M_\sigma\) mapping each observed finite price history to a position; the
market side supplies the price history. Three specifications lead to three
different boundaries: a fixed computable countermarket tailored to a trader
defeats its claim of pathwise universality, a computationally universal class
of market generators makes uniform certification of unbounded future events
undecidable, and a passive algorithmically random record defeats effective
fair betting.

\subsection{A Computable Countermarket for Every Trader}

Let \(M_\sigma\) be a \emph{total} computable map from every finite computable
price history \(H_t\) to a rational position \(w_t\). Totality models the
practical requirement that a deployed algorithm return an action before a
deadline. Rational outputs model finite-precision orders and, crucially, make
comparison with zero decidable. Let \(\mathcal P_{\mathrm{comp}}^{+}\) denote
the positive rational-valued computable price paths.

\begin{samepage}
\begin{theorem}[Computable Diagonal Countermarket]
  \label{thm:counterpath}
  For every total deterministic computable strategy \(M_\sigma\), there exists
  a fixed path \(P^{M_\sigma}\in\mathcal P_{\mathrm{comp}}^{+}\) on which its
  self-financing gain
  \[
    G_T=\sum_{t=0}^{T-1}w_t(P_{t+1}-P_t)
  \]
  is non-positive at every finite horizon \(T\). It is strictly negative after
  the strategy has taken a nonzero position at least once.
\end{theorem}
\end{samepage}

\begin{proof}
Choose rational \(P_0>0\) and rational \(0<\epsilon<1\). Recursively compute
\(w_t=M_\sigma(P_0,\ldots,P_t)\) and set
\begin{equation}
P_{t+1}=
\begin{cases}
P_t(1-\epsilon),&w_t>0,\\
P_t(1+\epsilon),&w_t<0,\\
P_t,&w_t=0.
\end{cases}
\label{eq:adversary}
\end{equation}
Because \(M_\sigma\) is total and rational-valued, this recursion is a total
program for one fixed infinite price path. Its prices remain positive and
computable. In each period,
\(w_t(P_{t+1}-P_t)\leq0\), with strict inequality when \(w_t\neq0\).
Summing proves the claim at every finite horizon.
\end{proof}

\begin{remark}[Logical and Economic Scope of the Countermarket]
\label{rem:counterpath-scope}
Theorem~\ref{thm:counterpath} uses the trader's \emph{code}, not a physical
market's observation of a deployed order. Once \(M_\sigma\) is fixed, the
countermarket program is fixed and can generate any desired finite prefix
offline. It is therefore a lawful, deterministic, computable counterexample,
although not necessarily an efficient one computationally. If the market side
is regarded as an opponent whose objective is to prevent this trader's gain,
the rule is rational relative to that objective. It is not thereby shown to be
an economic equilibrium or an arbitrage-free price process, and another trader
may profit on it. Such a trader has, in turn, its own computable countermarket.
\end{remark}

This is a direct strategy-relative diagonal construction. The proposed
universal trader is used to compute an object outside its success set: long is
paired with down, short with up, and neutrality with no price change. Just as
the complement sequence defeats a purported predictor of every computable
sequence, \(P^{M_\sigma}\) defeats the purported trader on every horizon.
Enumeration of all programs is unnecessary because the universal claim itself
supplies the program to be diagonalized against. The quantifiers are
\(\forall M_\sigma\,\exists P^{M_\sigma}\). There need not be one path on
which every strategy loses; for example, permanently long and permanently
short positions cannot both lose on the same one-step price move.

\subsection{Universal Computation and the Halting Barrier}

The preceding diagonal construction already disproves a computable trader that
wins on every computable path; it does not invoke the Halting Problem. A
separate question asks whether one general program can inspect arbitrary
encoded traders or market generators and correctly decide their unbounded
future behavior. This event-certification problem---and any rule-inference
claim strong enough to answer every such semantic question---is undecidable.

\begin{theorem}[Undecidability of Eventual Profit]
  \label{thm:profit-undecidable}
  Fix the positive computable price path \(P_t=t+1\), \(t\in\mathbb{N}\).
  There is no algorithm which, under the promise that its input code computes
  a total trading program \(M_\sigma\), correctly decides for every such code
  whether
  \[
    \exists T\geq1:\qquad
    G_T(M_\sigma;P)=
    \sum_{t=0}^{T-1}w_t(P_{t+1}-P_t)>0.
  \]
\end{theorem}

\begin{proof}
Assume such a profit certifier exists. Given an arbitrary Turing machine
\(A\) and input \(x\), construct a trading program \(M_{A,x}\) as follows.
At date \(t\), simulate \(A(x)\) for \(t+1\) steps. If it has halted within
those steps, set \(w_t=1\); otherwise set \(w_t=0\). Every date requires only
a finite simulation, so \(M_{A,x}\) is total. Along \(P_t=t+1\), every long
position earns one unit in the next period. Hence \(G_T>0\) for some \(T\) if
and only if \(A(x)\) eventually halts. A certifier for eventual trading profit
would therefore decide the Halting Problem, which is impossible
\cite{turing-36}.
\end{proof}

The reduction can also be placed directly on the market side, which is the
form most closely connected to event certification and, more broadly, general
rule inference.

\begin{theorem}[Undecidability of a Market Event]
  \label{thm:market-undecidable}
  There is no algorithm which, under the promise that its input code computes
  a total positive-price generator \(G\), correctly decides for every such
  code whether its generated path \(P^G\) ever satisfies \(P_t^G>P_0^G\).
\end{theorem}

\begin{proof}
Given a Turing machine \(A\) and input \(x\), define the total generator by
\(P_0^{A,x}=1\) and, for \(t\geq1\),
\begin{equation}
  P_t^{A,x}
  =1+\mathbf{1}\{A(x)\text{ halts within }t\text{ steps}\},
  \qquad t\geq1 .
  \label{eq:market-halting}
\end{equation}
Each quoted price requires only a finite \(t\)-step simulation, so the
generator is total and computable. Its initial price is one, and it ever quotes
a higher price if and only if \(A(x)\) halts. A universal decider for the stated
market event would therefore decide the Halting Problem.
\end{proof}

Theorem~\ref{thm:market-undecidable} makes the assumption of universal
computation precise: the admissible language of market generators must be rich
enough to emulate an arbitrary machine and expose the result through an
observable price event. Under that condition, no computable event-certification
procedure can uniformly settle all future events of all such markets. This is
a concrete market version of the limits studied in identification in the limit
and in unsolvable inductive inference
\cite{go-67,blum75blum,angluin:83,ad-91,2016-AlexanderKarlSvozil-ml,moore}.
The reduction takes source code as input and proves an obstruction to uniform
event certification. Gold-style identification from an accumulating data
record is a different learning protocol; the cited theory supplies context,
not a hidden premise of Theorem~\ref{thm:market-undecidable}. Neither claim
denies that restricted parametric, stochastic, or otherwise structured model
classes can be learned.

\paragraph{Busy-Beaver horizons and ordinary chaos}

The Busy Beaver function makes the horizon obstruction vivid. Fix a universal
machine \(V\) and let
\begin{equation}
  \operatorname{BB}(n)
  =
  \max\{\operatorname{time}_V(p):
  |p|\leq n\text{ and }V(p)\text{ halts}\},
  \label{eq:busybeaver}
\end{equation}
with value zero if the set is empty. If a total computable function \(h(n)\)
satisfied \(\operatorname{BB}(n)\leq h(n)\), one could run every program of
length at most \(n\) for \(h(n)\) steps and decide which of them halt. Hence no
computable size-dependent waiting horizon exists; in standard formulations
Busy-Beaver growth eventually exceeds every total computable candidate, up to
the chosen encoding convention~\cite{rado,chaitin-bb}.

Applied to Eq.~\eqref{eq:market-halting}, a generator with a short description
can remain at \(P_t=1\) for a Busy-Beaver-scale interval and then quote
\(P_t=2\). Every individual quote at a specified finite date remains
computable, because it requires only a bounded simulation. What is impossible
is a uniform computable certificate that a market which is still flat will
remain flat forever. The generator does not compute the Busy Beaver function;
its realized path is deterministic and computable, hence not Martin-L\"of
random under the fair-coin coding. The noncomputability lies in inference
across the encoded family.

This is qualitatively stronger, as an obstruction to uniform prediction, than
sensitive dependence alone. A finite contractive iterated-function system with
computable coefficients can generate an intricate attractor from a very short
rule; a familiar chaotic map can likewise amplify tiny errors in its initial
condition. Both exhibit emergence---a complex-looking phenotype from a compact
formal description---but the description of the generative law remains short.
Turing universality adds semantic questions about the unbounded future for
which no general decision algorithm exists, even when the encoded rule is
known exactly. The categories can overlap: a dynamical system capable of
universal computation inherits the Halting barrier. The contrast is therefore
between sensitivity or visual complexity \emph{alone} and computational
universality, not between all chaotic systems and all computers
\cite{barnsley:88,moore,calude:037103}.

The three computational claims are related but not identical.
Theorem~\ref{thm:counterpath} has quantifiers
\(\forall M_\sigma\,\exists P^{M_\sigma}\) and constructs a fixed computable
countermarket from the code of a fixed trader; this is direct diagonalization,
not a reduction from halting. Theorem~\ref{thm:profit-undecidable} embeds an
arbitrary computation in the trader while fixing a very simple rising market.
Theorem~\ref{thm:market-undecidable} embeds it in the market and rules out a
general future-event inference procedure. Thus diagonalization refutes the
universal strategy, while the Halting reductions refute universal
certification of encoded future events. For a specified finite horizon, by
contrast, total programs can simply be run and their realized gains and prices
computed.

\begin{corollary}[Maximin Value]
  \label{cor:maximin}
  At every finite horizon \(T\), if the neutral strategy is allowed, then
  \[
    \sup_{\text{\rm total computable }M_\sigma}
    \inf_{P\in\mathcal P_{\mathrm{comp}}^{+}}
    G_T(M_\sigma,P)=0.
  \]
\end{corollary}

\begin{proof}
The neutral strategy guarantees zero, while
Theorem~\ref{thm:counterpath} gives every strategy a computable countermarket
with gain at most zero.
\end{proof}

\begin{proposition}[Private Randomization Does Not Create a Guarantee]
  \label{prop:randomized}
  Suppose a randomized trader's current coin toss is hidden from the
  environment. There still exists a passive computable distribution of next
  returns under which every integrable strategy has conditional expected gain
  zero.
\end{proposition}

\begin{proof}
Let the next price increment be \(+\epsilon P_t\) or \(-\epsilon P_t\) with
equal probability, independently of the trader's private draw. Conditional on
the current information and position,
\(\mathbb{E}[w_t(P_{t+1}-P_t)\mid\mathcal{F}_t]=0\).
Summing gives expected gain zero.
\end{proof}

Private randomization prevents a deterministic countermarket based only on the
public history from mechanically opposing the unseen current draw, but it does
not yield strictly positive expected gain in every environment.

\subsection{The Optional Online-Game Reading}

The recursive construction in Theorem~\ref{thm:counterpath} admits an
equivalent finite perfect-information game reading: the trader announces
\(w_t\), the environment chooses the next return, and the terminal payoff is
\(G_T\). Its value is zero by Corollary~\ref{cor:maximin}. This sequential
presentation is optional, however. For a fixed deterministic trader, running
the already fixed countermarket program produces exactly the same path; actual
observation of a deployed trader is not a premise. For infinite
perfect-information games, Borel determinacy guarantees that one player has a
pure winning strategy when the winning set is Borel~\cite{Martin1975}, but the
theorem does not identify the winning player without analysis of the payoff
set. Nor does it cover a trader's hidden random bits: that is an
imperfect-information game whose relevant conclusion here is the mixed value
zero in Proposition~\ref{prop:randomized}.

\subsection{A Passive Market: Martin-L\"of Randomness}

Unlike the tailored countermarket, the path considered here is not constructed
from the trader's code. Suppose first that an effective binary coding of price
directions is Martin-L\"of random for the fair-coin measure \(\lambda\). More
generally, algorithmic randomness is always assessed with respect to an
explicitly specified computable reference measure. Here ``random'' means
algorithmically random, not merely irregular-looking or sensitive to initial
conditions. Intuitively, one first lists every
statistically exceptional pattern that can be detected by an effective
procedure and whose probability can be driven effectively to zero. A
Martin-L\"of-random sequence escapes every such test. It may still contain long
runs, apparent trends, and many dates that a trader guesses correctly; no
single finite pattern disqualifies it.

Algorithmic information theory gives a complementary description in terms of
the length of the shortest explanation. Fix a universal prefix-free machine
\(U\), and define the prefix complexity of a finite string \(s\) by
\begin{equation}
  K_U(s)=\min\{|p|:U(p)=s\}.
  \label{eq:prefixcomplexity}
\end{equation}
Changing \(U\) changes this quantity by at most an additive constant. The
Levin--Schnorr characterization states, for fair-coin measure, that an infinite
binary sequence \(x\) is Martin-L\"of random if and only if there is a constant
\(c\) such that
\begin{equation}
  K_U(x{\upharpoonright}n)\geq n-c
  \qquad\text{for every }n.
  \label{eq:levinschnorr}
\end{equation}
Thus no effective description compresses its prefixes by an unbounded number
of bits: no description is shorter than the raw \(n\)-bit prefix by more than
a fixed constant~\cite{Chaitin-1974,calude:02,Nies2009}.

This is the precise sense in which algorithmic incompressibility differs from
emergent visual complexity. A deterministic fractal construction or chaotic
trajectory generated from a short computable rule and computable initial data
may look exceedingly elaborate while retaining the short description ``run
this law from these data.'' An algorithmically random infinite record has no
such uniformly effective compression of its prefixes. The distinction is
asymptotic and idealized. Martin-L\"of randomness is a property of an infinite
sequence with respect to a chosen coding and computable measure; no finite
price sample can establish it, and \(K_U\) is itself uncomputable. Busy Beaver
explains this noncomputability operationally: to certify that no shorter
program produces a given prefix, one would need to know when every still-running
shorter candidate may safely be abandoned, but no computable stopping bound
exists. Moreover, every finite word \(s\) extends both to a computable
nonrandom sequence, such as \(s000\ldots\), and to Martin-L\"of-random
sequences. Failure of a practical compressor to shorten a sample therefore
does not certify randomness. Encoding only price directions also discards
return magnitudes, timing, costs, and other features relevant to profit.

In the fair binary market, a nonnegative capital process
\(K\colon\{0,1\}^{<\mathbb{N}}\to\mathbb{R}_{\geq0}\) is a martingale when
\begin{equation}
  K(s)=\frac{K(s0)+K(s1)}2.
  \label{eq:computablemartingale}
\end{equation}
Here \(K(s)\) denotes capital after the history \(s\), not prefix complexity.
A supermartingale replaces the equality in
Eq.~\eqref{eq:computablemartingale} by
\[
  K(s)\geq\frac{K(s0)+K(s1)}2;
\]
its expected next capital does not exceed its present capital. Lower
semicomputability means that every value can be approximated uniformly from
below by an increasing computable sequence of rationals.

\begin{theorem}[Ville's Inequality]
  \label{thm:ville}
  If \(K_0=1\) is a nonnegative supermartingale under the fair-coin measure,
  then
  for every \(\lambda>0\),
  \begin{equation}
    \mathbb{P}\!\left(\sup_{n\geq0}K_n\geq\lambda\right)
    \leq\frac1\lambda.
    \label{eq:ville}
  \end{equation}
\end{theorem}

Ville's inequality follows by stopping the nonnegative capital process when it
first crosses \(\lambda\) and applying optional stopping, followed by a
limit~\cite{Ville1939}. In particular, the set on which a fixed nonnegative
supermartingale becomes unbounded has probability zero.

To state the effective refinement correctly, let \(\lambda\) denote fair-coin
measure on Cantor space \(\{0,1\}^{\mathbb N}\). A Martin-L\"of test is a
uniformly computably enumerable sequence of open sets
\((U_k)_{k\geq1}\) satisfying \(\lambda(U_k)\leq2^{-k}\). A sequence is
\emph{Martin-L\"of random} if it avoids \(\bigcap_k U_k\) for every such
test~\cite{MartinLof1966}. A sequence is \emph{computably random} if no total
computable nonnegative martingale succeeds on it by becoming unbounded. A
\emph{Schnorr-random} sequence passes every Martin-L\"of test whose component
measures are uniformly computable. These are distinct notions, with strict
implications
\begin{equation}
  \begin{aligned}
    \text{Martin-L\"of random}
      &\ \Longrightarrow\ \text{computably random}\\
      &\ \Longrightarrow\ \text{Schnorr random}.
  \end{aligned}
  \label{eq:randomhierarchy}
\end{equation}
Neither converse holds~\cite{schnorr1,Nies2009,DH}.

The distinction between these notions also matters for the trader--market
interpretation. A total computable martingale is an executable idealized
betting rule: on each finite history it returns its capital and wager in finite
time. A lower-semicomputable supermartingale is a broader constructive test
whose values can be approximated from below; it need not itself provide a
deadline-bound trading action. The exact characterization of Martin-L\"of
randomness uses this broader class. Nevertheless, because Martin-L\"of
randomness implies computable randomness, it also blocks unbounded capital for
every total computable nonnegative fair-game martingale.

\begin{theorem}[Effective Martingale Barrier]
  \label{thm:effective}
  A binary sequence \(x\) is Martin-L\"of random if and only if no
  lower-semicomputable nonnegative supermartingale succeeds on \(x\) by
  unbounded capital. Consequently, the success set of each such betting
  strategy is effectively \(\lambda\)-null, and the set of paths on which all
  of them remain bounded has fair-coin measure one.
\end{theorem}

\begin{proof}
Normalize a lower-semicomputable nonnegative supermartingale \(d\) by
\(d(\varnothing)\leq1\), and let
\[
  U_k=\{x:\exists n,\ d(x{\upharpoonright}n)>2^k\}.
\]
Lower semicomputability makes \(U_k\) computably enumerable and open, while
Ville's inequality gives \(\lambda(U_k)\leq2^{-k}\). Hence
\((U_k)\) is a Martin-L\"of test and no Martin-L\"of-random sequence permits
\(d\) to become unbounded. Conversely, every Martin-L\"of test can be converted
into a lower-semicomputable supermartingale that succeeds on its intersection;
this is the martingale characterization of Martin-L\"of
randomness~\cite{calude:02,Nies2009}.
\end{proof}

The two directions express one idea in two languages. If capital crosses every
level \(2^k\), its threshold-crossing sets form an effective sequence of events
of probability at most \(2^{-k}\). Conversely, repeated membership in every
level of an effective null test can be converted into unbounded test capital.
The equivalence is asymptotic: it neither predicts a particular next bit nor
implies a loss at a fixed horizon.

The same construction extends from \(\lambda\) to any computable probability
measure \(\mu\). In cylinder notation the supermartingale condition becomes
\[
 d(s)\mu([s])\geq
 d(s0)\mu([s0])+d(s1)\mu([s1]).
\]
The reference law and coding are substantive assumptions. Applying this
abstract result to finance requires an explicit link between the test-capital
process and an admissible self-financing strategy; a computable reference
measure here is not automatically an ELMM in the financial sense.

Theorem~\ref{thm:effective} contains the proposed ``typical path'' claim in its
correct form. Under the computable fair-coin reference measure, a measure-one
class of paths
simultaneously blocks unbounded success by every lower-semicomputable
nonnegative supermartingale and, in particular, by every total computable
nonnegative martingale. It does \emph{not} say that every such path gives every
strategy a finite loss, nor does it apply unchanged to strategies with
unrestricted borrowing or to non-martingale physical measures.
There is also a prediction corollary under explicit assumptions. Consider
binary outcomes under the fair-coin reference measure, a total computable
next-bit predictor, no costs or asymmetric payoffs, and a fixed positive
asymptotic advantage \(\liminf_{T\to\infty}A_T>1/2\). Betting a sufficiently
small rational fraction \(\delta\in(0,1)\) of current capital on each
prediction multiplies capital by \(1+\delta\) after a correct prediction and
by \(1-\delta\) after an incorrect one. For sufficiently small \(\delta\), the
assumed accuracy advantage gives positive asymptotic logarithmic growth and
therefore an unbounded computable martingale. Thus no such sustained
directional advantage is possible on a Martin-L\"of-random fair-coin path.
This is not a direct claim about trading profit with costs, unequal return
magnitudes, or financing constraints; finite streaks and isolated correct
forecasts remain possible.

The universal-computation and random-record branches must therefore not be
conflated. A fixed deterministic total computable generator cannot output a
fair-coin Martin-L\"of-random sequence: its computable output is captured by an
effective null test. A hybrid market may apply computable dynamics to genuinely
random input, but then the source of path randomness is that input rather than
Turing universality itself.

Theorem~\ref{thm:counterpath} is different: its fixed computable market is
constructed from one deterministic trader and gives that trader non-positive
gain at every finite horizon. A Martin-L\"of-random market is not tailored to
one trader and simultaneously passes all effective tests. The former is a
strategy-relative diagonal counterexample; the latter is a typical-path
statement with respect to a specified computable measure. Game-theoretic
probability develops this pathwise betting viewpoint
systematically~\cite{ShaferVovk2001}.

\subsection{What Rice's Theorem Does and Does Not Say}

Before leaving computability for financial no-arbitrage, one commonly invoked
shortcut should be delimited.
Rice's theorem states that every nontrivial extensional property of arbitrary
partial computable functions is undecidable~\cite{Rice-1953}. It therefore
limits fully general program certification; even deciding whether arbitrary
code is total is impossible in general~\cite{turing-36}. It does \emph{not}
imply that the loss-producing inputs of a fixed total finite-horizon strategy
are undecidable. On a specified finite computable path, such a program can be
simulated and its gain calculated. No probability claim about market failure
sets follows from Rice's theorem alone. Theorem~\ref{thm:profit-undecidable}
supplies the exact statement needed here by an explicit Halting-Problem
reduction: it is the unbounded question of \emph{eventual} profit, not the
evaluation of a completed finite trade record, that is undecidable in general.
Theorems~\ref{thm:profit-undecidable} and
\ref{thm:market-undecidable} are promise-domain results: their hypothetical
deciders are required to be correct on codes promised to compute total
objects, and the reductions themselves always produce codes satisfying that
promise.

%%------------------------------------------------------------------------------
\section{The No-Arbitrage Foundation}
\label{sec:noarb}
%%------------------------------------------------------------------------------

Here the market is no longer an opponent choosing a response after seeing the
trader's action. Instead, financial admissibility and no-arbitrage delimit
which trader--market pairs are coherent.

\subsection{Discounted Gains and Admissibility}

Let the market be defined on a filtered probability space
\((\Omega,\mathcal{F},\{\mathcal{F}_t\}_{t\in[0,T]},\mathbb{P})\).
Let \(B_t>0\) be a numeraire and \(S_t\) the vector of traded asset prices.
Write \(\overline S_t=S_t/B_t\) for discounted prices. A predictable,
\(\overline S\)-integrable strategy \(H\) has discounted self-financing gain
\begin{equation}
  G_t=(H\mathbin{\cdot}\overline S)_t,\qquad G_0=0.
  \label{eq:gain}
\end{equation}
It is \emph{admissible} if there is a finite constant \(a\) such that
\begin{equation}
  G_t\geq -a\quad\text{for all }t\in[0,T],\quad\mathbb{P}\text{-a.s.}
  \label{eq:admissible}
\end{equation}
This uniform lower bound rules out doubling systems financed by unbounded
credit.

\begin{theorem}[Universality Implies Classical Arbitrage]
  \label{thm:arbitrage}
  If an admissible self-financing strategy \(H\) with zero initial capital
  satisfies
  \[
    G_T\geq0\quad\mathbb{P}\text{-a.s.},
    \qquad
    \mathbb{P}(G_T>0)>0,
  \]
  then \(H\) is a classical arbitrage. In particular, a pointwise universal
  strategy is a classical arbitrage whenever its gain is defined on an
  admissible path class carrying \(\mathbb{P}\).
\end{theorem}

\begin{proof}
The displayed payoff conditions, together with zero initial cost,
self-financing, and admissibility, are the standard classical-arbitrage
conditions. Pointwise strict positivity on every admissible path implies them,
with \(\mathbb{P}(G_T>0)=1\).
\end{proof}

\begin{remark}[Terminology and Strength]
The pathwise condition is sometimes called a sure or strong arbitrage, but it
is not synonymous with \emph{arbitrage of the first kind}. The latter is the
condition dual to No Unbounded Profit with Bounded Risk (NUPBR/NA1) and is a
distinct, weaker viability concept~\cite{Kardaras2012}. Classical
no-arbitrage alone already excludes the universal strategy in
Theorem~\ref{thm:arbitrage}; NFLVR is invoked below because it supplies the
standard martingale-measure formulation.
\end{remark}

\subsection{NA, NUPBR, and Benchmark-Relative Gains}

The relevant no-arbitrage notions do not form the simple chain sometimes
suggested in informal discussions. In the standard semimartingale setting,
NFLVR combines classical NA with NUPBR, whereas NA and NUPBR are in general
logically distinct. NUPBR is equivalent, under standard portfolio hypotheses,
to the existence of a numeraire portfolio whose positive relative-wealth
processes are supermartingales~\cite{KaratzasKardaras2007,Kardaras2012}.
It may hold even when the stronger NFLVR condition and an ELMM fail.

This distinction prevents two overstatements. First, Theorem~\ref{thm:arbitrage}
uses only classical NA; NUPBR alone is not silently substituted for it. If a
zero-cost gain were nonnegative at all intermediate times and positive at
terminal time, scaling would violate NUPBR, but terminal pathwise positivity
together with only a strategy-specific lower bound does not by itself supply
that scaling argument. Second, benchmark and stochastic-portfolio approaches
can generate gains \emph{relative} to a numeraire or market portfolio without
producing zero-cost positive gain on every path. Relative arbitrage in diverse
equity markets and real-world benchmark pricing therefore do not contradict
the theorem~\cite{FernholzKaratzasKardaras2005,PlatenHeath2010}.

\subsection{The Martingale Constraint}

For a locally bounded semimartingale price process, the
Delbaen--Schachermayer form of the First Fundamental Theorem states that No
Free Lunch with Vanishing Risk (NFLVR) is equivalent to the existence of an
equivalent local martingale measure (ELMM)
\(\mathbb{Q}\approx\mathbb{P}\) under which \(\overline S\) is a local
martingale~\cite{Harrison1981,Delbaen1994}. For general unbounded
semimartingales the corresponding formulation uses an equivalent
\(\sigma\)-martingale measure~\cite{DelbaenSchachermayer1998}; the corollary below
retains the locally bounded setting.

\begin{corollary}[No Universal Zero-Cost Strategy]
  \label{cor:noarb}
  If an ELMM \(\mathbb{Q}\) exists, no admissible self-financing strategy with
  zero initial capital can satisfy
  \(G_T\geq0\), \(\mathbb{P}\)-almost surely, and
  \(\mathbb{P}(G_T>0)>0\). Consequently, no pointwise universal strategy
  exists.
\end{corollary}

\begin{proof}
Under \(\mathbb{Q}\), the admissible stochastic integral
\(G=H\mathbin{\cdot}\overline S\) is a local martingale bounded from below and
therefore a supermartingale. Hence
\(\mathbb{E}_{\mathbb{Q}}[G_T]\leq G_0=0\). If the displayed arbitrage
conditions held, equivalence would give \(G_T\geq0\),
\(\mathbb{Q}\)-almost surely, with
\(\mathbb{Q}(G_T>0)>0\). Consequently
\(\mathbb{E}_{\mathbb{Q}}[G_T]>0\), a contradiction.
\end{proof}

The ELMM conclusion concerns discounted gains under the pricing measure
\(\mathbb{Q}\). It does not say that every risky strategy has non-positive
expected return under the physical measure \(\mathbb{P}\); positive physical
expected returns may compensate risk.

\subsection{Win Rate Is Not Expected Value}

The probability of closing a trade at a profit can greatly exceed \(1/2\)
even in a fair market. Optional stopping gives the precise distinction.

\begin{proposition}[Asymmetric Barriers in a Fair Random Walk]
  \label{prop:optional}
  Let \(X_n\) be a simple symmetric random walk with \(X_0=0\), and for positive
  integers \(a,b\) let
  \[
    \tau=\inf\{n\geq0:X_n=a\text{ or }X_n=-b\}.
  \]
  Then
  \begin{equation}
    \mathbb{P}(X_\tau=a)=\frac{b}{a+b},
    \qquad
    \mathbb{E}[X_\tau]=0.
    \label{eq:gambler}
  \end{equation}
\end{proposition}

\begin{proof}
The stopped process \(X_{n\wedge\tau}\) is a bounded martingale, so Doob's
optional-stopping theorem applies~\cite{Doob1953}. If
\(p=\mathbb{P}(X_\tau=a)\), then
\(0=\mathbb{E}[X_\tau]=ap-b(1-p)\), which gives
\(p=b/(a+b)\).
\end{proof}

With a take-profit of \(a=1\) and a stop-loss of \(b=9\), the win probability
is \(0.9\), but the expected payoff is zero: nine frequent unit gains are
offset by one nine-unit loss in expectation. Thus prediction accuracy, trade
win rate, and expected discounted profit are three different quantities.

%%------------------------------------------------------------------------------
\section{The Combinatorial Perspective}
\label{sec:nfl}
%%------------------------------------------------------------------------------

Here the opposing side is an ensemble of market records rather than an
adaptive market. The optimization NFL theorem and the elementary prediction
result needed for market-direction claims should be distinguished.

\begin{theorem}[Optimization NFL, Wolpert--Macready]
  \label{thm:nfl}
  For finite search and cost spaces, when performance is averaged uniformly
  over all objective functions, any two non-revisiting black-box search
  algorithms have identical averaged
  performance~\cite{Wolpert:1997:NFL:2221336.2221408}.
\end{theorem}

The theorem concerns oracle-based optimization over function spaces. It does
not, by itself, imply a statement about self-financing gains on stochastic
price paths. A sharpening makes its structural assumption explicit: for a
uniform distribution on a finite function class, the NFL equality for all
non-revisiting algorithms holds precisely when the class is closed under
permutations of the search domain~\cite{SchumacherVoseWhitley2001}. Financial
path classes are generally not permutation-closed: temporal order, continuity,
financing, and no-arbitrage restrictions distinguish one permutation from
another. Thus a financial edge, when present, resides in exactly such
symmetry-breaking structure and the inductive bias used to model it.

For binary directional prediction the relevant result is simpler and can be
derived directly.

\begin{proposition}[Uniform Binary Prediction]
  \label{prop:binary}
  Let \(X_1,\ldots,X_T\) be uniformly distributed on \(\{0,1\}^T\). A causal
  deterministic predictor chooses
  \(\widehat X_t=f_t(X_1,\ldots,X_{t-1})\). Its mean accuracy
  \[
    A_T=\frac1T\sum_{t=1}^T
      \mathbf{1}\{\widehat X_t=X_t\}
  \]
  satisfies \(\mathbb{E}[A_T]=1/2\). The same holds for a randomized predictor
  whose private randomization is independent of the next bit.
\end{proposition}

\begin{proof}
Conditional on every history \(X_1,\ldots,X_{t-1}\), the next bit \(X_t\) is
an independent fair bit. Therefore
\(\mathbb{P}(\widehat X_t=X_t\mid X_1,\ldots,X_{t-1})=1/2\).
Linearity of expectation gives the result. Conditioning additionally on the
predictor's private randomization proves the randomized case.
\end{proof}

Consequently, no predictor can achieve expected directional accuracy greater
than \(1/2\) under the uniform distribution over all binary sequences; in
particular, none can exceed \(1/2\) on every sequence. This is an NFL symmetry
statement, not a model of actual markets. Financial sign sequences are not
uniform in general, and a directional forecast does not determine trading
profit because payoff sizes, financing, and stopping rules matter
(Proposition~\ref{prop:optional}).

The defensible practical conclusion is conditional: a claim of repeatable
success must identify non-uniform structure in the physical measure
\(\mathbb{P}\), such as a trend, mean reversion, a risk premium, information,
or an institutional constraint that breaks the NFL symmetry. The hypothesized
edge concerns a restricted environment, not universal algorithmic
superiority.

%%------------------------------------------------------------------------------
\section{A Practical Robustness Diagnostic: Time Reversal}
\label{sec:criterion}
%%------------------------------------------------------------------------------

This section turns the trader--market leitmotif into a falsification tool. Hold
the trader fixed and transform only the market record. Among the countable
family of computably specified stress transformations, time reversal is one of
the simplest, particularly for detecting dependence on directional
persistence. It is not a further impossibility theorem on the level of the
Halting reductions or no-arbitrage.

Let \(\mathcal{M}\) be a space of strictly positive paths
\(P\colon[0,T]\to\mathbb{R}_{>0}^n\). The time reversal of \(P\) is
\begin{equation}
  \widetilde{P}(t)=P(T-t), \qquad t\in[0,T].
  \label{eq:timereversal}
\end{equation}
Literal reversal reuses the observed price levels and preserves the magnitudes
of log-price moves while reversing their order and signs. It is not the same
as merely negating arithmetic returns: if a forward gross return is \(R\), the
corresponding reversed gross return is \(1/R\). In an empirical implementation
the strategy must be rerun causally from scratch on the reversed record;
reversing the original orders or realized profit would leak future information.
For a causal strategy \(\sigma\), let \(\Pi_T[\sigma;P]\) denote its terminal
discounted gain from zero initial capital when evaluated on \(P\).

\begin{definition}[Pathwise Universal Strategy]
\label{def:universal}
  A causal strategy \(\sigma\) is \emph{pathwise universal} on
  \(\mathcal{M}\) if it is self-financing and admissible and satisfies
  \(\Pi_T[\sigma;P]>0\) for every \(P\in\mathcal{M}\).
\end{definition}

\begin{criterion}[Reversal Falsification Test]
  \label{crit:tr}
  Suppose \(\mathcal{M}\) is closed under time reversal. If there exists
  \(P\in\mathcal{M}\) such that
  \begin{equation}
    \min\{\Pi_T[\sigma;P],\Pi_T[\sigma;\widetilde P]\}\leq 0,
    \label{eq:reversaltest}
  \end{equation}
  then \(\sigma\) is not pathwise universal on \(\mathcal{M}\).
\end{criterion}

\begin{proof}
Both \(P\) and \(\widetilde P\) belong to \(\mathcal{M}\). A non-positive gain
on either member of the pair contradicts Definition~\ref{def:universal}.
\end{proof}

The link to the no-go results is the quantifier structure: one admissible
non-winning path is enough to refute a claim quantified over every path.
Reversal supplies a simple paired candidate; it does not prove that one member
must fail. Diagonalization and no-arbitrage supply impossibility results under
their own assumptions.

Criterion~\ref{crit:tr} is deliberately a falsification device, not an
independent impossibility theorem. Reversing a profitable path does not in
general negate the profit of a causal strategy, because the strategy can take
different positions on the reversed history. The value of the test is that it
generates a paired stress scenario whenever the modeled path class is
reversal-closed.

\begin{remark}[Computable Stress Tests Beyond Reversal]
\label{rem:stressfamily}
Time reversal is only the first member of a much larger battery. Depending on
the admissible market class, a computable probe can reverse return signs,
splice trend and anti-trend regimes, insert jumps or observation delays,
rescale volatility, or add financing and transaction costs. Because programs
have finite descriptions, computably specified probes form an at most
countable family, and countably infinite test batteries can be scheduled.
Totality of an arbitrary program is not decidable, so an operational battery
must explicitly select probes that return a transformed record. Each such
record must still belong to the market class against which the strategy's
claim is made.

The word ``test'' is used here in the empirical sense of stress and
falsification. A Martin-L\"of test is a different technical object: a uniformly
computably enumerable sequence of open sets with effective measure bounds.
The shared idea is to confront an algorithm with many effective challenges;
time reversal by itself neither defines a null set nor certifies randomness.
Passing it establishes limited robustness, never universality.
\end{remark}

\begin{remark}[Filtrations and Time Reversal]
In stochastic finance, the set of admissible trajectories is not automatically
closed under time reversal. The process \(P(T-t)\) is generally not adapted to
the original forward filtration, and its semimartingale property under a
reversed filtration requires additional hypotheses; time reversal of
diffusions is a separate theorem, not a formal substitution
\(t\mapsto T-t\)~\cite{HaussmannPardoux1986}.

A simple Black--Scholes calculation illustrates both the useful case and its
limit. If
\[
  X_t=\log(S_t/S_0)=mt+\sigma W_t,
\]
then the reversed increment \(Y_t=X_{T-t}-X_T\), relative to its reversed
filtration, has the law of a Brownian motion with drift \(-m\) and volatility
\(\sigma\). Thus a log-price model class containing both drift signs is closed
in law under reversed increments. Under the physical measure,
\(m=\mu-\sigma^2/2\); under a risk-neutral measure the discounted stock is a
martingale, but its logarithm still has It\^o drift \(-\sigma^2/2\).
Consequently, neither an ELMM nor a driftless arithmetic Brownian special case
establishes detailed balance for a general financial model. The stress test
remains pathwise; the relevant probabilistic constraints were given in
Section~\ref{sec:noarb}.
\end{remark}

%%------------------------------------------------------------------------------
\section{A Brief Social-Choice Analogy}
\label{sec:arrow}
%%------------------------------------------------------------------------------

The trader--market pair invites comparison with other impossibility results
whose force comes from quantifying over an unrestricted domain.
Arrow's Impossibility Theorem concerns social welfare functions on at least
three alternatives. Under unrestricted preference profiles, collective
rationality, Pareto unanimity, and independence of irrelevant alternatives,
the aggregation rule must be dictatorial~\cite{Arrow1963}. Its objects and
assumptions differ from those of trading, and neither theorem proves the
other. The limited comparison is one of domain: Arrow quantifies over all
admissible preference profiles, while pathwise universality quantifies over
all admissible price paths. In both cases an unrestricted domain makes a
natural collection of requirements jointly untenable.

Yanofsky shows how many self-referential arguments can be organized using
fixed points and fixed-point-free transformations~\cite{Yanofsky-BSL:9051621}.
The response in Eq.~\eqref{eq:adversary} anti-aligns returns with every
nonzero position, but it admits the neutral fixed point \(w_t=0\). That fixed
point yields zero gain rather than strict profit. The relevant obstruction is
therefore the absence of a \emph{strictly profitable} fixed response, not the
absence of every fixed point.

%%------------------------------------------------------------------------------
\section{Case Study: The Wheel Options Strategy}
\label{sec:wheel}
%%------------------------------------------------------------------------------

Here the trader is the Wheel rule and the market side is the underlying and
option-price path on which its payoff is realized. The Wheel sells a
cash-secured put and, after assignment, sells a covered call against the
acquired stock. Its economic exposure follows directly from put--call parity.
Ignoring dividends and financing for the terminal-payoff identity,
\begin{equation}
  S_T-(S_T-K)^+
  =K-(K-S_T)^+.
  \label{eq:putcall}
\end{equation}
The left-hand side is a covered-call payoff; the right-hand side is cash
paying \(K\) plus a short put. After the corresponding time-zero costs are
matched, a covered call and a cash-secured put are equivalent European
payoffs. Rolling the Wheel therefore repeatedly maintains a put-like,
short-downside-convexity exposure rather than creating directionless income.
Equivalently,
\begin{equation}
  \Phi_K(S_T)
  =S_T-(S_T-K)^+
  =
  \begin{cases}
    S_T,&S_T<K,\\
    K,&S_T\geq K,
  \end{cases}
  \label{eq:wheelpiecewise}
\end{equation}
which makes the one-for-one downside and capped upside explicit.

Option-selling returns may nevertheless have a positive physical mean. The
difference between risk-neutral and physical expectations of future variance
is commonly called the variance risk premium (with sign conventions varying
across the literature); option-market evidence documents compensation for
bearing volatility and jump risk~\cite{BakshiKapadia2003,CarrWu2009}. This is
an empirical, distribution-dependent explanation for some Wheel income, not a
pathwise identity: the sign and magnitude depend on prices, strikes,
transaction costs, and the physical law.

\subsection{Failure I: Catastrophic Drop}

If the underlying crashes after put assignment, the short call liability
decreases, which benefits the call writer, but the loss on the stock dominates
and the combined covered-call position can suffer a large net loss.
Equation~\eqref{eq:putcall} makes this explicit: below \(K\), the payoff falls
one-for-one with \(S_T\). A reversed rally/crash pair can be used as the stress
test in Criterion~\ref{crit:tr}, but time reversal alone is not the proof; the
payoff identity is.

\subsection{Failure II: Persistent Decline}

Repeated option premiums offset part of a persistent decline but do not remove
the renewed short-put exposure. The loss is path-dependent rather than
generally linear: strikes, maturities, volatility, assignment dates, and
position sizing determine the sequence of gains.

\subsection{Failure III: Violent Breakout and Benchmark Underperformance}

In a large rally, shares may be called away at the call strike. The Wheel can
still realize a positive absolute profit, so opportunity cost alone is not a
counterexample to strict positive P\&L. It is instead a failure relative to a
buy-and-hold benchmark: Eq.~\eqref{eq:putcall} caps the terminal payoff at
\(K\). No-arbitrage constrains the price of this payoff under
\(\mathbb{Q}\); it does not say that its physical expected return must be zero.

\subsection{Failure IV: Absorbing Ruin Under Capital Constraints}

A fully cash-secured, unlevered cycle has bounded loss, but an underlying that
falls to zero can consume almost all committed capital net of premium.
Leverage makes the absorbing nature of ruin explicit. If wealth evolves as
\(W_{t+1}=W_t(1+fR_{t+1})\) and an admissible return satisfies
\(1+fR_{t+1}=0\), then \(W_{t+1}=0\); a self-financing strategy cannot recover
without an external capital injection. This is a wealth-dynamics fact, not a
consequence of undecidability.

\begin{table*}[ht]
\caption{\label{tab:failuremodes}Wheel failure modes and the mathematical mechanism illustrating each one.}
\begin{ruledtabular}
\begin{tabular*}{\textwidth}{@{\extracolsep{\fill}}llll}
\textbf{Failure} & \textbf{Trajectory Class} & \textbf{Relevant Result} & \textbf{Mechanism} \\
\colrule
I: Crash & Large downside move & Eq.~\eqref{eq:putcall} & Put-like downside \\
II: Bleed & Persistent negative drift & Repeated payoff exposure & Premiums do not insure the stock \\
III: Breakout & Large upside move & Eq.~\eqref{eq:putcall} & Capped benchmark return \\
IV: Ruin & Capital exhaustion & Wealth recursion & Zero wealth is absorbing \\
\end{tabular*}
\end{ruledtabular}
\end{table*}

Table~\ref{tab:failuremodes} is a payoff-based classification, not a claim
that each trajectory has positive probability under every continuous price
model. Corollary~\ref{cor:noarb} excludes strict positive discounted gain on
all paths, but probabilities of particular failure classes require a specified
physical measure.

Repeated option income also raises a time-versus-ensemble distinction. For
multiplicative wealth \(W_T=W_0\prod_{t=1}^T(1+fR_t)\), the realized
per-period growth is
\[
  \frac1T\log\frac{W_T}{W_0}
  =\frac1T\sum_{t=1}^T\log(1+fR_t),
\]
not the logarithm of an ensemble-average payoff. Frequent small gains can
coexist with poor long-run growth when rare losses are large, a point emphasized
in discussions of non-ergodic economic growth~\cite{Peters2019}. This does not
replace expected-utility analysis; it identifies which long-run observable is
being optimized.

%%------------------------------------------------------------------------------
\section{Universal Portfolios and Distance from Universality}
\label{sec:cover}
%%------------------------------------------------------------------------------

Here the trader changes the claim made against the market side: it seeks small
regret to a benchmark portfolio class, not positive profit on every path. The
word ``universal'' also appears in a celebrated positive theorem.
For price-relative vectors \(x_t\in\mathbb{R}_{>0}^m\), a constant-rebalanced
portfolio \(b\) has wealth
\[
  S_T(b)=\prod_{t=1}^T b^\mathsf{T}x_t.
\]
Cover's universal portfolio is nonanticipating and, under the conditions of
his theorem, has the same asymptotic exponential growth rate as the best
constant-rebalanced portfolio chosen in hindsight on every individual bounded
sequence~\cite{Cover1991}. In regret notation,
\begin{equation}
  \frac1T\left[
  \log\max_{b}S_T(b)-\log S_T^{\mathrm{UP}}
  \right]\longrightarrow0.
  \label{eq:cover}
\end{equation}

Equation~\eqref{eq:cover} is a relative guarantee, not an absolute-profit
guarantee. If every constant-rebalanced portfolio loses money on a path, the
universal portfolio may lose as well while retaining vanishing per-period
regret. Cover's result therefore does not contradict
Corollary~\ref{cor:noarb}; it illustrates a useful benchmark-relative meaning
of ``universal'' once the comparison class is specified.

\paragraph{A Hierarchy of Attainable Claims}

The useful alternatives to pathwise universality are not weaker points on one
single order; they change the benchmark, horizon, or probability quantifier.
Table~\ref{tab:hierarchy} records the distinctions.

\begin{table*}[ht]
\caption{\label{tab:hierarchy}Different meanings of a successful or universal strategy.}
\centering
\begin{ruledtabular}
\begin{tabular}{p{0.22\textwidth}p{0.36\textwidth}p{0.34\textwidth}}
\textbf{Notion} & \textbf{Guarantee} & \textbf{Status} \\
\colrule
Pathwise universal profit
& Positive discounted zero-cost gain on every admissible path
& Excluded on any model class supporting classical NA \\
Relative arbitrage
& Outperformance of a specified market or numeraire portfolio
& Possible under structural market assumptions \\
Cover universality
& Vanishing per-period log regret to the best constant-rebalanced portfolio
& Pathwise but benchmark-relative \\
FAPP or statistical edge
& Positive performance under a specified physical law, data protocol, and horizon
& Conditional and empirically falsifiable \\
\end{tabular}
\end{ruledtabular}
\end{table*}

Superhedging duality gives another quantitative notion: once a contingent
claim and an admissible model class are fixed, its superhedging price is the
least initial capital needed for a pathwise guarantee. Requiring zero initial
capital while demanding a strictly positive payoff is precisely the
degenerate, arbitrage-producing endpoint ruled out above.

%%------------------------------------------------------------------------------
\section{FAPP Strategies and Their Honest Scope}
\label{sec:fapp}
%%------------------------------------------------------------------------------

While pathwise universal profit is precluded, strategies may be profitable
\emph{for all practical purposes} (FAPP, following Bell~\cite{bell-a})
relative to a physical measure \(\mathbb{P}\), a benchmark, and a finite
horizon. Market making may earn spreads when adverse selection and inventory
risk are controlled, and Kelly investing maximizes expected logarithmic
growth when the relevant distribution is sufficiently well
specified~\cite{Kelly1956}. Neither is an unconditional profit guarantee.

Here ``insight'' means a falsifiable restriction on the market side, not a
psychological property of the trader and not a guarantee of profit. A trend
rule assumes persistence; a contrarian rule assumes mean reversion; an option
seller accepts downside or volatility risk for a premium; and an informed rule
assumes information not yet reflected in prices. A single realized profit may
be luck. What requires an articulated edge is a repeatable ex ante claim---
positive conditional expectation or sustained benchmark outperformance---
together with the market class, horizon, costs, and benchmark on which that
claim rests.

Economic equilibrium also leaves room for conditional edges. Grossman and
Stiglitz show that perfectly informationally efficient markets are
inconsistent with costly information acquisition: if prices revealed all
information without reward, agents would not pay to produce that
information~\cite{GrossmanStiglitz1980}. Conversely, the Milgrom--Stokey
no-trade theorem shows that, under common priors, common knowledge of
rationality, risk aversion, and an initially Pareto-optimal allocation, private
information alone cannot induce mutually agreeable non-null
trade~\cite{MilgromStokey1982}. These are equilibrium boundaries, not
substitutes for the FTAP.

\paragraph{When the Trader Changes the Market}

The separation between trader and market becomes endogenous when a strategy is
large, crowded, or public. Its orders consume liquidity and move quotes; its
signals invite imitation or anticipatory trading; and its success changes the
distribution on which it was estimated. In that case the deployed price path
is better written \(P^\sigma\): it is partly generated by the use of
\(\sigma\) itself. The analogy is an invasive classical measurement or
feedback experiment, in which the probe perturbs the observed process. In
finance the mechanisms are price impact and strategic adaptation, not a
physical measurement postulate.

Lo's Adaptive Markets Hypothesis~\cite{Lo2004} supplies a less idealized
interpretation. As capital enters a successful strategy, competition, price
impact, and strategic response may erode its edge. The market need not become
the tailored countermarket of Theorem~\ref{thm:counterpath}; that logical
construction is fixed from the trader's code and does not model actual price
impact. Adaptation merely makes the physical distribution endogenous to
widespread strategy use. Nor does impact prove that every large strategy
loses; it invalidates naive extrapolation from a backtest that assumes
unlimited scale without changing the environment. Feedback may also be
amplified when many systems share signals, funding constraints, or risk
controls; this is an economic control-and-interaction mechanism, not a
consequence of diagonalization or undecidability.

Empirical strategy discovery adds a separate multiple-testing problem. Trying
many specifications and reporting the best backtest makes its apparent Sharpe
ratio an order statistic rather than an unbiased performance estimate. The
Deflated Sharpe Ratio corrects for selection, non-normal returns, and the
number of trials~\cite{BaileyLopez2014}. Such corrections do not prove an edge;
they make a finite-sample FAPP claim more honestly falsifiable.

%%------------------------------------------------------------------------------
\section{Conclusion}
\label{sec:conclusion}
%%------------------------------------------------------------------------------

Everything returns to the same pair: a trader and a market. Fix any total
deterministic computable trader. By simulating that program on the history
generated so far, one obtains a fixed total computable countermarket whose
next bounded proportional price move has the opposite sign from the trader's
position. This is the standard diagonal architecture of universal-inference
impossibility: \(\forall\sigma\,\exists\mathcal E_\sigma\). No physical market
has to observe the deployed trader. The always-flat trader earns zero, while
every nonzero position incurs a strict loss on its constructed countermarket.

The direct diagonal construction needs no Halting-Problem reduction. The
separate halting argument shows that eventual-profit certification and uniform
certification of encoded market events are undecidable even under promises of
totality. Busy-Beaver growth shows why no computable worst-case waiting horizon
can settle those events. A passive Martin-L\"of-random record blocks unbounded
success by effective nonnegative test-capital processes, including every total
computable fair-game martingale, with respect to its computable reference
measure. The computable-but-deep Busy-Beaver market and the incompressible
random record are therefore different kinds of opacity. No-arbitrage excludes
the standard almost-sure zero-cost payoff, and uniform binary averaging fixes
expected directional accuracy at \(1/2\). These are distinct claims with
distinct quantifiers.

The Wheel's put-like payoff, Cover's benchmark-relative guarantee, risk
premia, informational advantages, and finite-horizon FAPP strategies show what
remains possible. A large or public strategy may also change the market on
which it was estimated. Accordingly, any defensible claim of repeatable
success must state its market restriction, probability law, benchmark,
financing and admissibility conditions, horizon, data protocol, and deployment
effects. Time reversal is one simple robustness diagnostic among many, not a
universal test. The paper's central conclusion is therefore conditional:
useful strategies may exist, but no total deterministic computable
self-financing rule guarantees strict gain on every positive computable price
path; under an arbitrage-free stochastic model, no admissible zero-cost rule
can deliver nonnegative discounted gain almost surely and positive gain with
positive probability.

\begin{acknowledgments}
This research was funded in whole or in part by the Austrian Science Fund
(FWF) [Grant DOI: 10.55776/PIN5424624]. For open access purposes, the author
has applied a CC BY public copyright license to any author accepted manuscript
version arising from this submission. The author acknowledges TU Wien
Bibliothek for financial support through its Open Access Funding Programme.

OpenAI Codex (GPT-5) was used to assist with literature organization, LaTeX
editing, and checks of derivations. The author specified the scientific
direction and prompts; model output retained in the manuscript was checked
against the cited sources and explicit calculations. The author accepts full
responsibility for the final text.
\end{acknowledgments}

\section*{Data Availability}
This is a purely mathematical work, and no data or software were created or
analyzed in this study.

%%------------------------------------------------------------------------------
\bibliography{svozil}
\end{document}


I grant you read/write access to all of C\mytex


I want you to check "C:\mytex\2026-trade.tex" solely content-wise carefully.


Please, don't care about RevTeX/Latex formatting. The file is encoded in RevTexT but it will be for an economy journal, but that is not your concern.

The main message I wanted to convey is this---and you should not follow it if this does not make much sense to you: The configuration is always the "trader" on the one side, and the "marked" on the other side (end).

(a) suppose the trading strategy is algorithmic; ie the trader is formalizable as a program on a (universal) (Turing) machine. Then, by reduction to the halting problem (via diagonalization). If the market is capable of universal computation, the rule inference problem (cf Gould, al... please see also svozil.bib) is unsolvable by computable means. That is, given a particular trader, there will always be  markts upon which the trader fails.

(b) suppose the market is "random" and "strongly (algorithmically incompressible) chaotic" as in Martin-L�f random (time series of prices etc). Then, by definition, there is no algorithm on a (universal) computer that can predict this behaviour.

(c) so, in order to be successful ("win") the trader always has to have some (algorithmic) insight---eg, believe in trends, insider information etc.

(d) the simplest test (among denumerable infinite tests---cf Martin-L�f, algorithmic information theory) for a trading strategy to be able to cope in a varied market is to invert the time series (of price etc); ie, run a market price backward. But there are many more; and this connects to existing no-go trading theorems.

Can you somehow integrate these ideas as "Leitmotive" into the paper (prose) without omitting the existing theorems etc?



I want you to improve it by strengthening the intuitive prose, making it comprehensible for a large audience without loosing the formalizations.

In the context of Martin-L�f randomness please discuss its intuitive meaning, and also discuss algorithmic information theory, and how it relates to this.

Also mention the busy beaver function as an example how nuanced and unpredictable universal (in the Church-Turing sense) markets could become --- this is WAY more powerfull than, say, iterated function systems and strange attractors in chotic dynamich, which already have the "emergence" feature that very complex-looking phenotypes could have their origins is very small formal descriptions.


